\begin{ex}
Số gồm 9 chục và 9 đơn vị là số nào?
\shortans{99}
\loigiai{
Số đó là 99.
}
\end{ex}

\begin{ex}
Số liền sau của số 79 là số nào?
\shortans{80}
\loigiai{
Số liền sau của 79 là 80.
}
\end{ex}

\begin{ex}
Số liền trước của số 60 là số nào?
\shortans{59}
\loigiai{
Số liền trước của 60 là 59.
}
\end{ex}

\begin{ex}
Tính: $43 + 35$.
\shortans{78}
\loigiai{
Ta có $43 + 35 = 78$.
}
\end{ex}

\begin{ex}
Tính: $86 - 52$.
\shortans{34}
\loigiai{
Ta có $86 - 52 = 34$.
}
\end{ex}

\begin{ex}
Trong phép trừ $75 - 23 = 52$, số bị trừ là bao nhiêu?
\shortans{75}
\loigiai{
Số bị trừ là 75.
}
\end{ex}

\begin{ex}
Trong phép cộng $34 + 45 = 79$, tổng là bao nhiêu?
\shortans{79}
\loigiai{
Tổng là 79.
}
\end{ex}

\begin{ex}
Bạn Hà có 24 cái kẹo, bạn Nga có 11 cái kẹo. Bạn Hà nhiều hơn bạn Nga bao nhiêu cái kẹo?
\shortans{13}
\loigiai{
Nhiều hơn: $24 - 11 = 13$ cái kẹo.
}
\end{ex}

\begin{ex}
Đoạn thẳng MN dài 48 cm, đoạn thẳng PQ dài 35 cm. Đoạn MN dài hơn đoạn PQ bao nhiêu xăng-ti-mét?
\shortans{13}
\loigiai{
Dài hơn: $48 - 35 = 13$ cm.
}
\end{ex}

\begin{ex}
Điền số thích hợp: $41 < \dots < 43$.
\shortans{42}
\loigiai{
Số ở giữa là 42.
}
\end{ex}

\begin{ex}
Tính giá trị biểu thức: $54 + 23 - 31$.
\shortans{46}
\loigiai{
Ta có $54 + 23 = 77$, rồi $77 - 31 = 46$.
}
\end{ex}

\begin{ex}
Tính giá trị biểu thức: $89 - 46 + 15$.
\shortans{58}
\loigiai{
Ta có $89 - 46 = 43$, rồi $43 + 15 = 58$.
}
\end{ex}

\begin{ex}
Tìm số tròn chục lớn nhất có hai chữ số.
\shortans{90}
\loigiai{
Số đó là 90.
}
\end{ex}

\begin{ex}
Tìm số tròn chục bé nhất có hai chữ số.
\shortans{10}
\loigiai{
Số đó là 10.
}
\end{ex}

\begin{ex}
Hiệu của số tròn chục lớn nhất và số tròn chục bé nhất có hai chữ số là bao nhiêu?
\shortans{80}
\loigiai{
Hiệu: $90 - 10 = 80$.
}
\end{ex}

\begin{ex}
Một cửa hàng có 65 bao xi măng, đã bán được 32 bao. Cửa hàng còn lại bao nhiêu bao xi măng?
\shortans{33}
\loigiai{
Còn lại: $65 - 32 = 33$ bao.
}
\end{ex}

\begin{ex}
Thùng cam có 42 quả, thùng quýt có 36 quả. Cả hai thùng có tất cả bao nhiêu quả?
\shortans{78}
\loigiai{
Tổng số quả: $42 + 36 = 78$ quả.
}
\end{ex}

\begin{ex}
Tìm một số biết rằng số đó cộng với 25 thì được 68.
\shortans{43}
\loigiai{
Số đó là: $68 - 25 = 43$.
}
\end{ex}

\begin{ex}
Tìm một số biết rằng số đó trừ đi 34 thì được 45.
\shortans{79}
\loigiai{
Số đó là: $45 + 34 = 79$.
}
\end{ex}

\begin{ex}
Trên tia số, điểm biểu diễn số 45 cách điểm biểu diễn số 60 bao nhiêu đơn vị?
\shortans{15}
\loigiai{
Khoảng cách: $60 - 45 = 15$ đơn vị.
}
\end{ex}

\begin{ex}
Tìm tổng của số lớn nhất có một chữ số và số bé nhất có hai chữ số.
\shortans{19}
\loigiai{
Số lớn nhất có 1 chữ số là 9, số bé nhất có 2 chữ số là 10. Tổng: $9 + 10 = 19$.
}
\end{ex}

\begin{ex}
Hai số có hiệu bằng 24, nếu số trừ là 35 thì số bị trừ là bao nhiêu?
\shortans{59}
\loigiai{
Số bị trừ: $24 + 35 = 59$.
}
\end{ex}

\begin{ex}
Hai số có tổng bằng 97, nếu một số hạng là 54 thì số hạng kia là bao nhiêu?
\shortans{43}
\loigiai{
Số hạng kia: $97 - 54 = 43$.
}
\end{ex}

\begin{ex}
Bác An năm nay 58 tuổi, chú Bình 36 tuổi. Bác An nhiều hơn chú Bình bao nhiêu tuổi?
\shortans{22}
\loigiai{
Nhiều hơn: $58 - 36 = 22$ tuổi.
}
\end{ex}

\begin{ex}
Điền số thích hợp: $30 + 40 + \dots = 95$.
\shortans{25}
\loigiai{
Ta có $30 + 40 = 70$, vậy số cần điền là $95 - 70 = 25$.
}
\end{ex}
